% Input-ready method fragment; no document class or package declarations.
% The parent document needs standard AMS math support.
% Citation keys are listed in docs/RELATED_WORK_V02.md.
% No effectiveness results or automated API-synthesis claims are made here.

\section{Evidence-Driven Representation and Rule Refinement}
\label{sec:method}

We study a family of static satellite--ground scheduling problems under
different resource-constraint configurations. Each configuration is fixed
during a scheduling run. Offline interventions provide evidence about how
decisions should adapt across configurations; deployment uses one frozen
program that reads the current conflict graph. Our central distinction is
between a rule that misuses available information and a representation that
cannot express the required rankings.

\subsection{Scheduling Model and Fixed Execution Kernel}
\label{sec:method:model}

Let $V$ be a set of contact opportunities with nonnegative rewards $w_v$.
A legal configuration $\theta$ induces a conflict graph
$G_\theta=(V,E_\theta,w)$, where an edge denotes joint infeasibility of
two contacts under the modelled resource constraints. Contact selection is
the maximum-weight independent set problem
\begin{equation}
 \max_{I\subseteq V}\ w(I)=\sum_{v\in I}w_v
 \quad\text{subject to } I\text{ independent in }G_\theta.
 \label{eq:method:mwis}
\end{equation}
This reduction assumes that all hard constraints in the evaluated model
are captured by the graph and any explicit boundary commitments.
Constraints involving higher-order resource interactions require an
extended model rather than an unsupported pairwise-feasibility claim.

A boundary $B=(F,X)$ consists of fixed selected contacts $F$ and explicitly
excluded contacts $X$. We require $F$ to be independent and $F\cap X$ to
be empty. With $N_\theta(F)$ denoting the union of the neighbours of $F$,
the residual candidates are
\begin{equation}
 R_\theta(B)=V\setminus\bigl(F\cup X\cup N_\theta(F)\bigr).
\end{equation}
A program $P=(\phi,h)$ computes features $\phi(G_\theta,R,v)$ and a finite
score $h(\phi(G_\theta,R,v))$ for each $v\in R$. The kernel selects a
highest-scoring contact, adds it to the schedule, and removes its closed
neighbourhood from $R$. Ties use a fixed deterministic convention; a tie
does not satisfy a strict ranking requirement. The final schedule is
checked independently against the graph.

\paragraph{Kernel invariant.}
Starting from a feasible $F$, each selected contact belongs to the residual
set and is therefore compatible with every previously selected contact.
Induction preserves independence. At least the selected contact is removed
at each step, so the finite process terminates. This establishes feasibility
for the graph model, not optimality of the learned ordering.

\subsection{Certified Relations from Aligned Interventions}
\label{sec:method:evidence}

An aligned intervention changes admissible resource parameters from
$\theta_0$ to $\theta_1$, rebuilding conflict edges while preserving contact
identities, rewards, objective units, and the intended boundary semantics.
The competing actions $a$ and $b$ must be distinct, adjacent in both graphs,
and individually feasible under $B$ in both configurations. In particular,
we reject comparisons whose apparent preference change is caused by an
action becoming infeasible.

Let $\mathcal{I}_\theta$ be the independent sets of $G_\theta$.
For $d\in\{a,b\}$, define the full conditional completion value
\begin{equation}
 \begin{split}
 V_\theta(d\mid B)=\max\{w(I):\ &I\in\mathcal{I}_\theta,\\
 &F\cup\{d\}\subseteq I,\ I\cap X=\varnothing\}.
 \end{split}
 \label{eq:method:conditional}
\end{equation}
The maximization includes all remaining contacts, not only a sampled local
region. The certificate scope is conditional on $B$; changing a boundary
changes the quantity being certified. A checked feasible completion for
each action--configuration combination establishes that all four spaces
are nonempty.

A budgeted oracle returns sound intervals
\begin{equation}
 L_\theta(d\mid B)\leq V_\theta(d\mid B)
 \leq U_\theta(d\mid B).
 \label{eq:method:bounds}
\end{equation}
Lower bounds are values of independently checked full-graph schedules.
Upper bounds must bound the entire conditional completion space, for
example through a valid relaxation or the open frontier of branch and
bound. For an oracle region $H$, let $A=R_\theta(B\cup\{d\})$, where
$B\cup\{d\}$ means adding $d$ to $F$. If $U_H$ bounds the MWIS value of
$G_\theta[A\cap H]$, one valid, possibly loose upper bound is
\begin{equation}
 U_\theta(d\mid B)=w(F\cup\{d\})+U_H
       +\sum_{v\in A\setminus H}w_v.
 \label{eq:method:localupper}
\end{equation}
It relaxes the conflicts involving the outside region. A local feasible
solution must be completed and checked globally before its value is used
as a lower bound. Expanding $H$ may tighten evidence, but budget exhaustion
never turns an unresolved interval into a certified label. Integer or
outward-rounded objective arithmetic is used to preserve bound validity.

For an objective tolerance $\epsilon\geq0$, the inequalities
\begin{equation}
 \begin{split}
 L_{\theta_0}(a\mid B)&>U_{\theta_0}(b\mid B)+\epsilon,\\
 L_{\theta_1}(b\mid B)&>U_{\theta_1}(a\mid B)+\epsilon
 \end{split}
 \label{eq:method:reversal}
\end{equation}
certify a strict preference reversal. They imply the same strict ordering
of the respective conditional optima by
Equation~\ref{eq:method:bounds}. If the same action satisfies the strict
inequality in both configurations, we certify preference preservation.
Preservation concerns ordering, not equality of objective values. Overlapping
intervals remain unresolved; a rollout average or an incumbent alone is
insufficient to certify the opposite bound.

Each certified comparison induces a relational specification
$r=(x^+,x^-)$, where $x=(G_\theta,R_\theta(B),v)$ is a decision occurrence
and $x^+$ is preferred. The requirement is
\begin{equation}
 h(\phi(x^+))>h(\phi(x^-)).
 \label{eq:method:requirement}
\end{equation}
Implementation checks may impose a separate numerical score margin.
That margin is not the oracle's objective tolerance $\epsilon$.
The evidence certifies a comparison of optimal completions under a
specified boundary. It does not guarantee that every heuristic completion
following the preferred action dominates every completion following the
other action, or that greedy use of all such relations is globally optimal.

\subsection{Detecting Representation-Induced Contradictions}
\label{sec:method:alias}

Let $\mathcal{D}$ be the finite set of certified requirements and
$\mathcal{Z}$ its decision occurrences. The representation includes every
value observable by the scoring function. Define exact feature aliasing by
\begin{equation}
 x\sim_\phi y\quad\Longleftrightarrow\quad \phi(x)=\phi(y).
\end{equation}
The quotient requirement graph $Q_\phi$ has one vertex for each equivalence
class in $\mathcal{Z}/\!\sim_\phi$ and an edge
$[x^+]\rightarrow[x^-]$ for every requirement. Self-loops are retained.
Equality uses exact canonical feature values; clustering or rounded
similarity is not used as a proof of aliasing.

\paragraph{Proposition (finite-sample representability).}
There exists an unrestricted deterministic pointwise real-valued function
on the observed feature vectors satisfying all strict requirements if and
only if $Q_\phi$ is acyclic. Thus a cycle or self-loop proves that no rule
over this representation can satisfy all requirements.

\paragraph{Proof.}
Because a pointwise function assigns the same score to identical feature
vectors, every edge requires a strict decrease between its endpoint
scores. A directed cycle would imply
$s_1>s_2>\cdots>s_k>s_1$, a contradiction; a self-loop requires
$s_1>s_1$. Conversely, in a finite acyclic graph assign to each vertex the
length of its longest outgoing path. Every edge then has a larger score at
its source than at its destination. These scores define a function on the
observed vectors satisfying every requirement. They can be scaled to any
fixed finite positive margin.\hfill$\square$

The converse is an abstract existence statement: a lookup mapping on a
finite set is not necessarily expressible by the permitted rule language,
and gives no guarantee on unseen vectors. Therefore an acyclic quotient
does not certify success of concrete DSL synthesis. Approximate aliases
are useful diagnostics but cannot establish this impossibility result.
The proposition applies to the current pointwise scoring interface, not
to arbitrary policies allowed to inspect extra graph structure, action
identifiers, or history. Instance and action identifiers are excluded from
synthesized semantic features to prevent label lookup.

A cycle returns its requirements and the underlying aliased graph
occurrences as a structural witness. The refinement loop distinguishes
this obstruction from a candidate rule merely failing an expressible
requirement. It first audits intervention alignment and certificate
validity, then requests distinguishing features. This makes representation
expansion a response to identified missing distinctions rather than an
unconditional increase in feature count.

\subsection{Typed Structural Features and Offline Synthesis}
\label{sec:method:typed}

The feature language distinguishes individual nodes, node sets, edge sets,
and numbers. Its operations expose the root contact, the available set,
neighborhood queries, typed set algebra, induced edges, counts, and
aggregates of contact weights. Edge aggregates include sums of endpoint
minimum weights and endpoint weight products, with bounded numerical arithmetic.
For example, if
$C(v)=R\setminus(\{v\}\cup N_\theta(v))$, both
\begin{equation}
 f_1(v)=\sum_{u\in C(v)}w_u,\qquad
 f_2(v)=\sum_{\{u,z\}\in E_\theta[C(v)]}\min(w_u,w_z)
\end{equation}
are deployable graph statistics. Equal $f_1$ values can hide different
internal compatibility among surviving opportunities; $f_2$ can expose
some such differences. These aggregates are features, not automatically
valid completion bounds. In particular, shared edge endpoints can cause
multiple penalties for the same contact.

The offline synthesizer receives relational specifications, cycle or
failure witnesses, the typed library, and computation limits. It proposes
feature ASTs and a corresponding rule AST. Type checks, expression-size
limits, finite-output checks, and a closed execution interface reject
invalid proposals. Generated programs cannot invoke the oracle, external
services, arbitrary imports, or saved labels at deployment.
The LLM proposes distinctions and formulas; certificate checking and
program evaluation determine acceptance.

The budgeted loop retains earlier specifications, replays them after each
revision, and records unresolved queries as well as successes. Optional
acquisition can request comparisons on training instances where retained
programs disagree. A disagreement is a query priority, not a certified
preference; it enters $\mathcal{D}$ only after the bound checks succeed.
Neither feature proposal nor budgeted search is claimed to converge to
a consistent or optimal program.

\subsection{Joint Selection, Deployment, and Scope}
\label{sec:method:selection}

Candidate selection jointly considers specification consistency,
complete-schedule quality, and feature cost. Let $s(P)$ be the fraction of
certified training requirements satisfied, $q(P)$ the normalized reward on
a separate selection set, and $c(P)$ the mean deterministic feature work.
The prototype uses the declared selection utility
\begin{equation}
 \begin{split}
 J(P)&=q(P)-\lambda\bigl(c(P)/c_{\rm ref}-1\bigr),\\
 \text{subject to }&s(P)\geq\tau\text{ and valid interfaces},
 \end{split}
 \label{eq:method:utility}
\end{equation}
where $c_{\rm ref}>0$ is the work of the weight-only baseline. Selection
reward is normalized by the optimum on the small selection problems where
that optimum is computed offline. The targeted variant additionally rejects
an exact-alias quotient containing a cycle. Utility ties favour higher
schedule quality, lower feature cost, and then earlier candidates.
The pilot configuration uses $\tau=0.75$ and $\lambda=0.002$; these are
declared experimental settings, not theoretical guarantees. All remaining
violations are reported, so passing the gate is not described as full
specification satisfaction. No selection score establishes unseen-instance
consistency or schedule optimality.

The deterministic work counter includes primitive dispatches, set and
membership scans, weight reads, and arithmetic or aggregation work. It is
a declared cost proxy, not an exact CPU-instruction count. Actual feature
and scoring time and total deployment latency are also recorded.
Feature caching is valid only for the graph and residual state on which
the feature depends. Offline costs separately
record LLM calls, candidate proposals, oracle work, and selection evaluations;
zero online calls do not imply zero discovery cost.

After selection, both $\phi$ and $h$ are frozen. Held-out evaluation separates
contact instances and constraint configurations from synthesis and selection,
including tests of their joint shift. The same program responds to each
current graph within the fixed kernel; no online LLM or optimization-oracle
query is allowed. Adaptation therefore means changes in the frozen program's
ranking as its deployable structural inputs change.

\paragraph{Prototype status and limitations.}
The initial executable prototype uses assistant-proposed candidates and
replayable offline artifacts. This is assisted synthesis, not evidence of
an autonomous API evolution run. Large-graph certification can be expensive
or inconclusive; pairwise graph modelling has a stated scope; the typed
library may omit the required distinction; finite evidence can overfit;
and feature costs can grow with graph density. A publishable evaluation
must retain candidate provenance and measure these limitations rather than
equating successful small examples with effective general adaptation.
